\documentclass[11pt,a4paper]{article}
\usepackage[margin=27mm]{geometry}
\usepackage{fontspec,unicode-math,amsmath,amsthm,booktabs,longtable,hyperref}
\setmainfont{STIXGeneral}\setmathfont{STIX Two Math}\hypersetup{hidelinks}
\newtheorem{theorem}{Theorem}\newtheorem{lemma}[theorem]{Lemma}
\newtheorem{proposition}[theorem]{Proposition}
\newcommand{\PSL}{\operatorname{PSL}}\newcommand{\Aut}{\operatorname{Aut}}
\newcommand{\Cl}{\operatorname{Cl}}\newcommand{\Cen}{\operatorname{C}}
\title{The complete correction to the locally projective quotient census}
\author{Application to the Hartley--Leemans 2006 erratum}\date{}
\begin{document}\maketitle
The two-page erratum \cite{erratum} removes eight purported quotients from the original census \cite{source}. We recover its full deletion and numerical contract from the complete common theory \cite{common}. That common theory proves the universal marked group before computing the census; no theorem from the separate application to the original article is an unpriced premise here. Its entire five-page table companion is part of the common acquisition and records all surviving rows, both facet counts, all containment links and all normalizer quotients.

\section{The corrected result}
Let $\mathcal P$ be the universal regular polytope with dodecahedral facets and hemi-icosahedral vertex figures. In the common paper, Theorems 6 and 7 establish its exact marked group
\[
 W=[5,3,5]/\langle\!\langle(s_1s_2s_3)^5\rangle\!\rangle
   \cong J_1\times\PSL_2(19),\qquad |W|=600415200.
\]
We use the same indexed generators and the same semisparse quotient convention as that paper. In particular the subgroup must satisfy the endpoint condition for every conjugate, not merely for a sample of conjugates.

\begin{theorem}\label{erratum}
Delete exactly the original historical rows
\[
13,\ 70,\ 78,\ 90,\ 91,\ 95,\ 111,\ 127.
\]
There remain $137$ conjugacy classes of semisparse subgroups of $W$, hence $137$ quotient isomorphism classes, with $26$ maximal semisparse classes. Exactly $66$ quotients are section regular: $65$ retain the original local types and one has only hemi-dodecahedral facets. The other $71$ have mixed dodecahedral and hemi-dodecahedral facets, with all vertex figures hemi-icosahedral.
\end{theorem}
\begin{proof}
The common paper's Lemma 1 tests the complete forbidden union of ambient element classes, so no conjugate can be missed. Its Theorem 8 applies that test to every one of the $1262$ subgroup conjugacy classes of the actual marked group. The resulting $137$-vertex containment graph has $349$ edges and $26$ maxima. Proposition 2 gives the facet counts and Lemma 3 detects exactly the maximal admissible classes. The output separates $65$ rows with no hemi-facets, one with no ordinary facets, and $71$ with both. The fixed vertex figure and the two regular facet types identify the first two classes of local behavior as the section-regular quotients.

The deletion labels are not inferred from the difference of totals. The common proof evaluates the thirty literal source maximal-subgroup generating sets and gives actual forbidden-conjugate witnesses for rows $13,70,95,127$. It then identifies rows $78,111$ as the unique cyclic subgroup classes of orders $45,15$ in row $13$, and row $90$ as the unique cyclic order-$30$ child of row $70$. Row $91$ is the nonabelian order-$30$ child with normalizer quotient of order $12$; the other such child has quotient order $36$ and is not confused with it. For every deleted row the saved permutations satisfy
\[
 x\in N,\qquad q\in W,\qquad x^q\in H_0H_3\setminus\{1,\omega\},
 \quad \omega=(s_0s_1s_2)^5.
\]
These eight explicit failures and the complete bijection of all surviving historical rows with the computed graph prove the literal correction.
\end{proof}

\section{Full output and source qualifications}
The complete table companion to \cite{common} retains all $137$ subgroup orders, both facet counts, every one of the $349$ maximal-containment links, actual subgroup descriptions and all quotient automorphism groups. It is not a selected summary table. Full permutation generators additionally specify each embedded subgroup, and the eight witness files specify the actual forbidden conjugacies. The whole comparison file \texttt{source\_tables.json} retains all $145$ historical rows. These complete outputs and data belong to the declared common boundary; neither historical \texttt{g.txt} nor a supplied semisparse list is used to enumerate the groups.

The source describes a defect in its old implementation. That diagnosis is a historical report. The complete conjugacy-class test proves the correction without reconstructing the old bug or its timing. Bibliographic matter, regret and acknowledgements remain source context.

The two-page publisher erratum is the exact source edition here. Its associated tables are read in the accessible fifteen-page 2003 arXiv article, not in the unseen twelve-page journal article. The arXiv Table 1 label $58$ corresponds to full-table row $63$, whereas the full-table row $58$ is a different subgroup. Its rows $21$ and $27$ have the explicit subgroup-column discrepancies discussed in the common paper. Row $41$ uses a noncyclic elementary abelian group of order eight; rows $38$--$40$ retain distinct actions despite similar abbreviated structure strings. None of these arXiv-only qualifications is silently assigned to the journal edition or to the erratum. The exact marked representatives, rather than orders alone, retain the required actions.

The retained external inputs are the general string-C-group/polytope and semisparse quotient correspondences, Hartley's general endpoint criterion, classical local spherical and antipodal groups, and the named exact finite-group algorithms and data. The formerly retained particular universal-group theorem is now proved in the common paper. All six complete common programs are required, not sampled validation. There is no new Lean declaration, historical-priority claim or new conjecture answer.
\begin{thebibliography}{9}
\bibitem{common} \emph{Marked groups and complete locally projective quotient censuses}, common paper, complete five-page census companion and six whole required programs in the present edition.
\bibitem{source} M. I. Hartley and D. Leemans, \emph{Quotients of a Universal Locally Projective Polytope of type $\{5,3,5\}$}, arXiv:math/0307075v1, 5 July 2003, fifteen pages. Journal reference: Math. Z. \textbf{247} (2004), 663--674; the journal edition has not been collated here.
\bibitem{erratum} M. I. Hartley and D. Leemans, erratum, Math. Z. \textbf{253} (2006), 433--434, doi:10.1007/s00209-005-0872-z.
\end{thebibliography}
\end{document}
